\section{Deadlock detection and recovery}
\label{sec:deadlock}

\begin{lede}
How the system tells a genuine deadlock apart from ordinary queueing --- which
in dense lifelong traffic is most of what it sees --- and the seven
increasingly drastic things it tries once it is sure.
\end{lede}

\subsection{Three distinct stall signals}
\label{sec:deadlock-signals}

Three related but different counters live on \code{Robot}, and they are worth
disambiguating precisely because they feed different formulas:

\begin{center}\small
\begin{tabular}{@{}p{0.22\linewidth}p{0.42\linewidth}p{0.28\linewidth}@{}}
\toprule
\textbf{Counter} & \textbf{Meaning} & \textbf{Reset} \\
\midrule
\code{waiting\_time} & steps spent physically stationary while not yet at the waypoint & on any move, or on reaching the waypoint \\
\code{blocked\_time} & consecutive steps the PIBT recursion returned \textsc{invalid} & on any \textsc{valid} result \\
\code{no\_progress\_steps} & consecutive steps $\Dg{w}(x_i)$ did not strictly decrease & on any route-distance improvement \\
\bottomrule
\end{tabular}
\end{center}

\subsection{Progress signal}

\begin{equation}
  \Delta_i(t) = \Dg{w}(\pos{i}{t-1}) - \Dg{w}(\pos{i}{t}).
\end{equation}
\code{update\_progress} treats a robot as not stalled whenever it is idle,
already at its waypoint, or at distance 0; otherwise it increments
\code{no\_progress\_steps} unless the distance strictly decreased and both ends
were finite. Then
\begin{equation}
  \mathrm{is\_blocked}(i) \iff
    \code{no\_progress\_steps} \geq t_{\mathrm{blocked}}
    \ \vee\
    \code{blocked\_time} \geq t_{\mathrm{blocked}},
  \qquad t_{\mathrm{blocked}} = 10.
\end{equation}

\subsection{Repeated configurations and the wait-for graph}
\label{sec:deadlock-detection}

\paragraph{Repeated configuration.} With period $k = 4$, on the 16-entry
position history deque:
\begin{equation}
  \mathrm{repeated}(i) \iff
    \mathrm{history}[-k{:}] = \mathrm{history}[-2k : -k],
\end{equation}
that is, the robot's last $k$ positions exactly equal the $k$ before them.

\paragraph{Wait-for dependency graph.} An edge $i \rightarrow j$ exists exactly
when robot $i$'s \code{waiting\_for\_robot}, set during the
priority-inheritance step of \Cref{alg:pibt}, is robot $j$. \code{find\_cycles}
runs an iterative white/grey/black DFS to find every simple directed cycle. A
cycle here is a genuine circular wait --- $i$ waits for $j$, \dots, waits for
$i$ --- the textbook signature of deadlock.

\code{detect\_deadlocked\_groups} combines two signals: any cycle in the
wait-for graph containing at least one currently-blocked robot, and a
BFS-connected-component clustering of the \emph{remaining} blocked robots
within $\max(2, R_{\mathrm{local}})$ hops of each other. The second pass
catches deadlocks that do not present as a clean wait-for cycle --- several
robots each stuck against static geometry in the same area --- without needing a
global traversal. Group lifetimes are tracked by \code{frozenset} membership
keys across timesteps, which is what makes ``detected'', ``recovered'' and
``unrecovered'' distinguishable.

\begin{caveat}
\textbf{Corroboration.} With \param{require\_deadlock\_corroboration} (the
default), a group must show the progress signal \emph{and} at least one of a
wait-for cycle or a repeated configuration. The specification named three stall
signals and the implementation treated the weakest as sufficient --- which in
dense lifelong traffic describes an ordinary queue, not a deadlock. Recovery
therefore escalated on healthy robots, and levels 5--7 below cost far more than
they saved: $22$ tasks per 1000 timesteps against $134$ on
\variant{warehouse\_corridors}. \Cref{sec:contribution-corroboration} explains
why this stayed invisible for so long.
\end{caveat}

\subsection{Seven-level progressive recovery}
\label{sec:recovery-levels}

\code{recover\_from\_deadlock} escalates \textbf{one level per timestep} a
given group persists, applying the next level's remedy each time the group is
still detected:

\begin{enumerate}
\item \textbf{Recompute routes} --- from scratch for every robot in the group.
  The cheapest possible fix: perhaps the route was merely stale.
\item \textbf{Recompute affected aisles} --- force \code{lock\_until} $= t$ on
  the group's current and next aisles, allowing an immediate re-decision
  instead of waiting out the lock; reopen an affected aisle if it is empty,
  force it to \textsc{draining} if it is occupied and directional.
\item \textbf{Release stale reservations} held by the group.
\item \textbf{Raise blocked priorities} ---
  $p_i \mathrel{+}= k_b \cdot \code{no\_progress\_steps}$, an escalating boost
  on top of \Cref{eq:priority}.
\item \textbf{Allow temporary reverse movement} ---
  \code{allow\_reverse\_until} $= t + t_{\mathrm{blocked}}$, with a shorter
  grace period for ignoring direction entirely.
\item \textbf{Assign escape vertices} --- route each member to its nearest
  unoccupied, unclaimed parking or passing bay, falling back to any
  degree-$\geq 3$ vertex, and set its state to \textsc{recovery}, which also
  grants $P_{\mathrm{emergency}}$.
\item \textbf{Run the local fallback planner} --- the last resort: drop all
  directional constraints for the group for $t_{\mathrm{deadlock}}$ steps,
  release reservations, restore task-directed state for anyone mid-recovery.
\end{enumerate}

If a group is still detected after level 7, that remedy is reapplied each
subsequent timestep and the group is counted once, not repeatedly, as
unrecovered.

\paragraph{Why one level per timestep and not all seven per call.} Progress
from a remedy can only be \emph{observed} on the next timestep, because the
simulator is synchronous. Running all seven in one call would apply six of them
blind. ``Progressive'' recovery in a synchronous loop necessarily means one
escalation per detection cycle.
